\scanpage{202}
% archive sheet number: 201
% source lines 1--2
$\unsure{m|\mathfrak{m}}/m^2$ dual de $P(U)$ primitifs, \deleted{\unsure{espace de}} \unsure{de dim} finie \unsure{sur}
\[
=\unsure{\text{tange}|\text{tang.}|\text{tué}}\ \text{par le \unsure{noyau} de }F.
\]
% source lines 3--10
s'il est de dim finie, le \uline{socle} est fini, donc on a un objet \uline{artinien} [N.B. l'injectif \inserted{indécomp.} de la catégorie des \unsure{bigèbres|algèbres} de Dieudonné est artinien \unsure{...}], donc on a un module \deleted{\unsure{noeth.}} de Chevalley artinien = module de type. fini sur $A$. De plus, la condition est que $M/FM$ est de dim finie sur $k$ i.e. \uline{de longueur finie}.

% source centre-right enclosed conclusion
\noindent\fbox{\begin{minipage}{0.55\linewidth}
Ceci indique que le module $M$ est un module de torsion
\end{minipage}}
% source centre-left category diagram and its two exterior braces
\diagram{\begin{tikzcd}[ampersand replacement=\&,column sep=small]
A_F\ \left\{\vphantom{\begin{matrix}C\\C_1\\C_F\end{matrix}}\right.\&\&\unsure{C|\mathcal{C}}\&\&\left.\vphantom{\begin{matrix}C\\C_1\\C'_F\end{matrix}}\right\}\ A_S\&\text{\textcircled{\small $S=A-F$}}\\
\&\&C_1\arrow[u]\&\\
\&C_F\arrow[ur]\&\&C'_F\arrow[ul]\\
\&\&C_0\arrow[ul]\arrow[ur]\arrow[d,dash]\\
\&\&0
\end{tikzcd}}
% source lower-left note beside C_F/C_0
\margin{\unsure{cat. primaire} \deleted{\unsure{d'objets}} \unsure{engendrée} dans $C/C_0$ par l'objet $(A/FA)_{C_0}$ \uline{simple} $\unsure{\omega|w}$}
% source lower-middle note, linked by a brace to C'_F and C_0
\unsure{cat. primaire gauche}. \unsure{engendrée} dans $C_1/C_0$ par les autres objets simples de $C_1/C_0$, sur $F$ est \deleted{bijectif}\inserted{surjectif} dans l'objet mod $C_0$
\[
C_1/C_0\simeq C_F/C_0\times C'_F/C_0
\]
% source small drawing in the right margin; two crossing curved strokes
\figure{\begin{tikzpicture}[scale=0.7]
\draw[gray] (0,0) .. controls (0.8,1.1) and (1.6,1.8) .. (3,2);
\draw (1,2.1) .. controls (0.95,1.2) and (1.1,1) .. (1.65,0.55);
\fill (1.05,1.25) circle (1.3pt);
\end{tikzpicture}}
% source right-hand notes, over the printer's imprint
\margin{\uline{limites projectives d'objets artiniens}}
% source bottom-left bracketed fragments; the manuscript leaves the opening brackets unclosed
[objet \deleted{\unsure{simple}} \unsure{engendré} par $A/FA$

[objet \deleted{\unsure{noeth.}} \unsure{engendré;} par $A/VA$

\uline{\unsure{deux} objets \deleted{\unsure{non}} simples}
\[
\deleted{\unsure{A}}\unsure{F_2|F^2}
\]
\margin{\uline{limites inductives d'objets noethériens}}
% letterhead: upside down in lower right, overlapped by notes
\begin{flushright}
M. DUFOUR\\
VARITYPISTE\\
131, RUE LECOURBE\\
PARIS XV\textsuperscript{e}\\
$\blacksquare$\\
C. C. P. Paris 14.143-34\\
R. M. 140.029
\end{flushright}
